\subsection{点分布与函数的卷积}

设 G 为李群，X 为 \(C^{r}\) 类流形，\((g, x) \mapsto gx\) 为 G 在 X 上的一个 \(C^{r}\) 类左作用律。对所有 \(x \in X\)，用 \(\rho(x)\) 表示 x 的轨道映射。

定义 3. 设 \(t \in \mathcal{T}^{(s)}(\mathrm{G})\)，\(s \leqslant r\)。设 \(f: \mathrm{X} \to \mathrm{F}\) 为取值于一个分离多范空间的 \(\mathbf{C}^r\) 类函数（例如 \(\mathrm{F} = \mathrm{K}\)）。\(t\) 与 \(f\) 的卷积，记作 \(t * f\)，是 \(\mathbf{X}\) 上以 \(\mathrm{F}\) 为值域、由

\[
(t * f) (x) = \langle t ^ {\vee} * \varepsilon_ {x}, f \rangle .
\]

于是

\[
\begin{array}{l l} (t * f) (x) = \langle \rho (x) _ {*} (t ^ {\vee}), f \rangle & (\text {no. 3, Proposition 14 (ii)}) \\ = \langle t ^ {\vee}, f \circ \rho (x) \rangle & (Diff. \& Anal. Man., R, 13.2.3) \\ = \langle t, (f \circ \rho (x)) ^ {\vee} \rangle & (Diff. \& Anal. Man., R, 13.2.3). \end{array}\tag{4}
\]

注意，定义 3 也可以写成更对称的形式

\[
\langle \varepsilon_ {x}, t * f \rangle = \langle t ^ {\vee} * \varepsilon_ {x}, f \rangle .\tag{5}
\]

函数 \((g,x)\mapsto f(gx) = (f\circ \rho (x))(g)\)（\(\mathbf{G}\times \mathbf{X}\) 上的）是 \(\mathbf{C}^r\) 类的。由 Differentiable and Analytic Manifolds, R, 13.4.4，函数 \(x\mapsto \langle t^\vee,f\circ \rho (x)\rangle\) 因而是 \(\mathbf{C}^{r - s}\) 类的（若 \(s < \infty\)）。换句话说，若 \(s < \infty\)，则 \(t*f\) 是 \(\mathbf{C}^{r - s}\) 类的。

显然，\(t * f\) 线性地依赖于 \(t\) 和 \(f\)。

特别地，公式 (4) 蕴含：对 \(g \in G\)，

\[
\left(\varepsilon_ {g} * f\right) (x) = f \left(g ^ {- 1} x\right)\tag{6}
\]

也就是

\begin{enumerate}
\def\labelenumi{(\arabic{enumi})}
\setcounter{enumi}{6}
\tightlist
\item
\end{enumerate}

\[
\varepsilon_ {g} * f = \gamma (g) f.
\]

设 \(\mathbf{K} = \mathbf{R}\) 或 \(\mathbf{C}\)，\(\mathbf{G}\) 与 \(\mathbf{X}\) 都是有限维的，且 \(\mathbf{X}\) 具有一个在 \(\mathbf{G}\) 下不变的正测度。\(\varepsilon_g * f\) 的定义与 Integration, Chapter VIII, § 4, no. 1 中的定义一致（参看~同处公式 (2)）。

命题 17. 设 \(t \in \mathcal{T}^{(s)}(\mathrm{G})\)，\(t' \in \mathcal{T}^{(s')}(\mathrm{X})\)，\(f: X \to \mathrm{F}\) 为 \(\mathbf{C}^{r}\) 类函数，满足 \(s + s' \leqslant r\)。则

\[
\langle t ^ {\prime}, t * f \rangle = \langle t ^ {\vee} * t ^ {\prime}, f \rangle .
\]

\[
\begin{array}{l l} \langle t ^ {\prime}, t * f \rangle = \langle t ^ {\prime}, x \mapsto \langle t, g \mapsto f (g ^ {- 1} x) \rangle \rangle & \text { by   (4) } \\ = \langle t \otimes t ^ {\prime}, (g, x) \mapsto f (g ^ {- 1} x) \rangle & \text {(Diff. \& Anal. Man., R, 13.4.4)} \\ = \langle t ^ {\vee} \otimes t ^ {\prime}, (g, x) \mapsto f (g x) \rangle & \text {(Diff. \& Anal. Man., R, 13.2.3)} \\ = \langle t ^ {\vee} * t ^ {\prime}, f \rangle . \end{array}
\]

命题 18. 设 \(t \in \mathcal{T}^{(s)}(\mathrm{G})\)，\(t' \in \mathcal{T}^{(s')}(\mathrm{G})\)，\(f: \mathrm{X} \to \mathrm{F}\) 为 \(\mathbf{C}^r\) 类函数，满足 \(s + s' \leqslant r\)。则

\[
(t * t ^ {\prime}) * f = t * (t ^ {\prime} * f).
\]

对所有 \(x \in X\)，

\[
\begin{array}{l l} \langle \varepsilon_ {x}, (t * t ^ {\prime}) * f \rangle = \langle (t * t ^ {\prime}) ^ {\vee} * \varepsilon_ {x}, f \rangle & \text {by (5)} \\ = \langle t ^ {\prime \vee} * (t ^ {\vee} * \varepsilon_ {x}), f \rangle & \text {(Propositions 2 and 7)} \\ = \langle t ^ {\vee} * \varepsilon_ {x}, t ^ {\prime} * f \rangle & \text {(Proposition 17)} \\ = \langle \varepsilon_ {x}, t * (t ^ {\prime} * f) \rangle & \text {(Proposition 17).} \end{array}
\]

若 \(r \geqslant \infty\)，则我们看到，\(C^\infty\) 类函数（\(X\) 上、以 \(F\) 为值域）的集合是代数 \(\mathcal{T}^{(\infty)}(G)\) 上的一个左模。

命题 19. 设 \(t \in \mathcal{T}^{(s)}(\mathrm{G})\)，\(s \leqslant r\)。设 \(f\)（相应地 \(f'\)）为 \(\mathrm{C}^r\) 类函数，定义在 \(\mathbf{X}\) 上、取值于一个分离多范空间 \(\mathrm{F}\)（相应地 \(\mathrm{F}'\)）。设 \((u, u') \mapsto uu'\) 为 \(\mathrm{F} \times \mathrm{F}'\) 到一个分离多范空间 \(\mathrm{F}''\) 中的连续双线性映射，使得 \(ff'\) 是 \(\mathrm{C}^r\) 类函数，定义在 \(\mathbf{X}\) 上、取值于 \(\mathrm{F}''\)。设 \(\sum_{i=1}^{n} t_i \otimes t_i'\) 为 \(t\) 在余积下在 \(\mathcal{T}^{(s)}(\mathrm{G}) \otimes \mathcal{T}^{(s)}(\mathrm{G})\) 中的像。则

\[
t * (f f ^ {\prime}) = \sum_ {i = 1} ^ {n} (t _ {i} * f) (t _ {i} ^ {\prime} * f ^ {\prime}).
\]

设 \(x \in X\)，并用 \(\rho(x)\) 表示 x 的轨道映射。于是

\[
\begin{array}{l l} \langle \varepsilon_ {x}, t * (f f ^ {\prime}) \rangle = \langle t ^ {\vee}, (f f ^ {\prime}) \circ \rho (x) \rangle & \text {by (4)} \\ = \langle t ^ {\vee}, (f \circ \rho (x)) (f ^ {\prime} \circ \rho (x)) \rangle \\ = \sum_ {i = 1} ^ {n} \langle t _ {i} ^ {\vee}, f \circ \rho (x) \rangle \langle t _ {i} ^ {\prime \vee}, f ^ {\prime} \circ \rho (x) \rangle \\ & (\text {Diff. \& Man. Anal., R, 13.5.2}) \\ = \sum_ {i = 1} ^ {n} \langle \varepsilon_ {x}, t _ {i} * f \rangle \langle \varepsilon_ {x}, t _ {i} ^ {\prime} * f ^ {\prime} \rangle & \text {by (4).} \end{array}
\]

注 1. 设 G 为李群，X 为 \(\mathbf{C}^r\) 类流形，\((x,g)\mapsto xg\) 为 G 在 X 上的一个 \(\mathbf{C}^r\) 类右作用律。若 \(t\in \mathcal{T}^{(s)}(\mathrm{G})\)（\(s\leqslant r\)）且 \(f\colon \mathrm{X}\to \mathrm{F}\) 是 X 上的一个 \(\mathbf{C}^r\) 类函数，我们用 \(f*t\) 表示 X 上由

\[
\begin{array}{r l} \langle \varepsilon_ {x}, f * t \rangle & = \langle \varepsilon_ {x} * t ^ {\vee}, f \rangle \\ & = \langle \rho (x) _ {*} (t ^ {\vee}), f \rangle \\ & = \langle t ^ {\vee}, f \circ \rho (x) \rangle \\ & = \langle t, (f \circ \rho (x)) ^ {\vee} \rangle . \end{array}\tag{8}
\]

特别地

\[
(f * \varepsilon_ {g}) (x) = f (x g ^ {- 1})\tag{9}
\]

也就是

\[
f * \varepsilon_ {g} = \delta (g) ^ {- 1} f.\tag{10}
\]

在明显的记号下，命题 17、18、19 成为

\begin{enumerate}
\def\labelenumi{(\arabic{enumi})}
\setcounter{enumi}{10}
\tightlist
\item
\end{enumerate}

\[
\langle t ^ {\prime}, f * t \rangle = \langle t ^ {\prime} * t ^ {\vee}, f \rangle\tag{11}
\]

\[
f * (t * t ^ {\prime}) = (f * t) * t ^ {\prime}\tag{12}
\]

\[
(f f ^ {\prime}) * t = \sum_ {i = 1} ^ {n} (f * t _ {i}) (f ^ {\prime} * t _ {i} ^ {\prime}).\tag{13}
\]

命题 20. 设 \(G\)、\(G'\) 为李群，\(X\) 为 \(C^r\) 类流形，\((g, x) \mapsto gx\)（相应地 \((x, g') \mapsto xg'\)）为 \(C^r\) 类的、\(G\)（相应地 \(G'\)）在 \(X\) 上的左（相应地右）作用律。设 \((gx)g' = g(xg')\) 对所有 \(x \in X\)、\(g \in G\)、\(g' \in G'\) 成立。设 \(t \in \mathcal{T}^{(s)}(G)\)、\(t' \in \mathcal{T}^{(s')}(\mathrm{G}')\) 且 \(f: X \to F\) 为 \(C^r\) 类函数，满足 \(s + s' \leqslant r\)。则

\[
(t * f) * t ^ {\prime} = t * (f * t ^ {\prime}).
\]

对所有 \(x \in X\)，

\[
\begin{array}{l l} \langle \varepsilon_ {x}, (t * f) * t ^ {\prime} \rangle = \langle \varepsilon_ {x} * t ^ {\prime \vee}, t * f \rangle & \text { by   (8) } \\ = \langle t ^ {\vee} * (\varepsilon_ {x} * t ^ {\prime \vee}), f \rangle & \text {(Proposition 17)} \\ = \langle t ^ {\vee} * \varepsilon_ {x}, f * t ^ {\prime} \rangle & \text {(Proposition 2 and (11))} \\ = \langle \varepsilon_ {x}, t * (f * t ^ {\prime}) \rangle & \text { by   (5). } \end{array}
\]

特别地，把 G 看作通过左平移和右平移作用在它自身上。若 \(f \colon \mathbf{G} \to \mathbf{F}\) 是 G 上的 \(\mathbf{C}^r\) 类函数且 \(t \in \mathcal{T}^{(s)}(\mathbf{G})\)（\(s \leqslant r\)），则 \(t * f\) 与 \(f * t\) 当 \(s < \infty\) 时是 G 上的 \(\mathbf{C}^{r - s}\) 类函数。进而，设 \(t' \in \mathcal{T}^{(s')}(\mathbf{G})\) 满足 \(s + s' \leqslant r\)。则

\[
(t * f) * t ^ {\prime} = t * (f * t ^ {\prime}).\tag{14}
\]

特别地，\(\mathcal{C}^{\infty}(\mathrm{G})\) 是一个 \((\mathcal{T}^{(\infty)}(\mathrm{G}),\mathcal{T}^{(\infty)}(\mathrm{G}))\)-双模。公式 (5) 与 (8) 的特殊情形是

\[
\langle t, f \rangle = \left\langle \varepsilon_ {e}, t ^ {\vee} * f \right\rangle = \left\langle \varepsilon_ {e}, f * t ^ {\vee} \right\rangle .\tag{15}
\]

注 2. 设 \((g,x)\mapsto gx\) 为一个 \(\mathbf{C}^r\) 类的、\(G\) 在 \(X\) 上的左作用律。设 \(t\in \mathbf{U}_s(\mathbf{G})\)（\(s\leqslant r\)），\(\Omega\) 为 \(X\) 的开子集，\(f\colon \Omega \to F\) 为 \(\mathbf{C}^r\) 类函数。\(t*f\) 也可以由公式 (4) 或 (5) 来定义；它是定义在 \(\Omega\) 上、取值于 \(F\) 的函数，是 \(\mathbf{C}^{r - s}\) 类的（若 \(s < \infty\)）。本节的结果以明显的方式推广到这一情形。

